% Rapport de projet 5MLDE : chaîne MLOps de prévision des ventes (UCI Online Retail)
% Compilation : latexmk -pdf main.tex   (ou pdflatex main.tex deux fois)
\documentclass[11pt,a4paper]{article}

% ------------------------------------------------------------------ à compléter
\newcommand{\auteur}{Nema Elisée Kourouma \\[2pt] Ju-heon Hwang \\[2pt] Ahmed Chaudhry}
\newcommand{\auteurcourt}{Kourouma, Hwang, Chaudhry}
\newcommand{\formation}{5MLDE -- Machine Learning \& Data Engineering}
\newcommand{\etablissement}{Nom de l'établissement}  % <- ton école
\newcommand{\depot}{https://github.com/Elie224/5MLDE}
\newcommand{\datedoc}{Novembre 2026}

% ------------------------------------------------------------------ paquets
\usepackage[T1]{fontenc}
\usepackage[utf8]{inputenc}
\usepackage{lmodern}
\usepackage[french]{babel}
\usepackage[margin=2.3cm]{geometry}
\usepackage{microtype}
\usepackage{amsmath}
\usepackage{graphicx}
\usepackage{booktabs}
\usepackage{tabularx}
\usepackage{array}
\usepackage{multirow}
\usepackage{enumitem}
\usepackage[table]{xcolor}
\usepackage{tikz}
\usetikzlibrary{positioning,arrows.meta,fit,backgrounds,shapes.geometric}
\usepackage{pgfplots}
\pgfplotsset{compat=1.18}
\usepackage{listings}
\usepackage{tcolorbox}
\tcbuselibrary{skins,breakable}
\usepackage{fancyhdr}
\usepackage{titlesec}
\usepackage[hidelinks]{hyperref}
\usepackage{caption}
\usepackage[section]{placeins}

% ------------------------------------------------------------------ style
\definecolor{bleu}{HTML}{1F4E79}
\definecolor{bleuclair}{HTML}{DCE6F1}
\definecolor{vert}{HTML}{2E7D32}
\definecolor{orange}{HTML}{E07B00}
\definecolor{rouge}{HTML}{C62828}
\definecolor{gris}{HTML}{5F6368}
\definecolor{fondcode}{HTML}{F5F7FA}

\titleformat{\section}{\Large\bfseries\color{bleu}}{\thesection}{0.8em}{}
\titleformat{\subsection}{\large\bfseries\color{bleu}}{\thesubsection}{0.7em}{}
\titleformat{\subsubsection}{\normalsize\bfseries}{\thesubsubsection}{0.6em}{}

\pagestyle{fancy}
\fancyhf{}
\fancyhead[L]{\small\color{gris}Chaîne MLOps -- prévision des ventes}
\fancyhead[R]{\small\color{gris}\auteurcourt}
\fancyfoot[C]{\small\thepage}
\renewcommand{\headrulewidth}{0.3pt}

\setlist{itemsep=2pt,topsep=4pt}
\setlength{\parskip}{4pt}
\setlength{\emergencystretch}{2.5em}
\renewcommand{\arraystretch}{1.25}
\captionsetup{font=small,labelfont={bf,color=bleu}}

\newcommand{\code}[1]{\texttt{\detokenize{#1}}}

\lstset{
  basicstyle=\ttfamily\small,
  backgroundcolor=\color{fondcode},
  frame=single, rulecolor=\color{bleuclair},
  breaklines=true, columns=fullflexible,
  xleftmargin=4pt, xrightmargin=4pt, aboveskip=6pt, belowskip=6pt,
  literate={é}{{\'e}}1 {è}{{\`e}}1 {à}{{\`a}}1 {ê}{{\^e}}1 {ô}{{\^o}}1 {→}{{$\rightarrow$}}1
}

\newtcolorbox{encadre}[1][]{enhanced,breakable,colback=bleuclair!40,colframe=bleu,
  boxrule=0.6pt,arc=2pt,left=6pt,right=6pt,top=4pt,bottom=4pt,fonttitle=\bfseries,#1}

% Capture d'écran : insère l'image si elle existe, sinon un cadre « à insérer ».
\newcommand{\capture}[3][0.92\linewidth]{%
  \begin{figure}[htbp]\centering
  \IfFileExists{../captures/#2}{\includegraphics[width=#1]{../captures/#2}}{%
    \fcolorbox{gris}{fondcode}{\parbox[c][4.2cm][c]{0.85\linewidth}{\centering\color{gris}
      Capture à insérer : \code{docs/captures/#2}\\[2pt]\small #3}}}%
  \caption{#3}\label{fig:#2}
  \end{figure}}

% ================================================================== document
\begin{document}

% ------------------------------------------------------------------ titre
\begin{titlepage}
  \centering
  \vspace*{1.5cm}
  {\large\color{gris}\formation\par}
  \vspace{0.3cm}
  {\large\color{gris}\etablissement\par}
  \vspace{2.2cm}
  {\color{bleu}\rule{\linewidth}{1pt}\par}
  \vspace{0.5cm}
  {\Huge\bfseries\color{bleu} Une chaîne MLOps de bout en bout\par}
  \vspace{0.4cm}
  {\LARGE Prévision des ventes hebdomadaires par produit\par}
  \vspace{0.3cm}
  {\large\itshape Cas UCI Online Retail, rejoué par lots hebdomadaires\par}
  \vspace{0.5cm}
  {\color{bleu}\rule{\linewidth}{1pt}\par}
  \vspace{2cm}
  {\Large \auteur\par}
  \vspace{0.4cm}
  {\large Rapport de projet -- travail en équipe\par}
  \vfill
  \begin{tabular}{rl}
    Dépôt & \url{\depot}\\
    Données & UCI Online Retail -- CC BY 4.0 -- DOI 10.24432/C5BW33\\
  \end{tabular}
  \vspace{1cm}

  {\large\datedoc\par}
\end{titlepage}

% ------------------------------------------------------------------ résumé
\section*{Résumé}
\addcontentsline{toc}{section}{Résumé}

Un modèle d'apprentissage automatique ne crée de la valeur qu'une fois en service, et il
ne la conserve que si l'on sait d'où viennent ses données, quelle version répond, comment
elle se comporte et comment revenir en arrière. Ce projet construit, autour d'un cas
métier volontairement simple, la chaîne qui apporte ces garanties.

Le cas retenu est le réassort d'un commerce en ligne britannique : prévoir, pour ses 50
produits les plus vendus, la quantité qui sera vendue la semaine suivante. Le jeu de
données historique \emph{UCI Online Retail} est rejoué semaine par semaine afin de simuler
l'arrivée de nouvelles ventes. Chaque lot est d'abord soumis à des contrôles de qualité
\textbf{bloquants} ; il alimente ensuite un entraînement dont le résultat n'est mis en
service que s'il fait au moins aussi bien que le modèle en place sur les mêmes semaines de
test. L'enchaînement est automatisé par Prefect, les modèles sont tracés et versionnés dans
MLflow, servis par une API FastAPI qui indique toujours la version qui répond, et toutes
les prédictions, erreurs et ventes réelles sont conservées dans PostgreSQL. Prometheus et
Grafana assurent la supervision.

La chaîne a été démontrée de bout en bout : sept versions du modèle ont été produites et
promues automatiquement, un lot invalide a été arrêté avant tout entraînement, une coupure
de la base a été absorbée sans interrompre le service et a déclenché l'alerte prévue, et un
retour arrière a été effectué en une commande. Les tests de charge conduisent à une
recommandation chiffrée : une unité de \textbf{2 workers, 1 vCPU et 1 Go} sert
\textbf{55,8 requêtes par seconde} avec une latence p95 de 110 ms et aucune erreur.

\vspace{0.3cm}
\noindent\textbf{Mots-clés :} MLOps, qualité des données, traçabilité, registre de modèles,
orchestration, supervision, dimensionnement.

\newpage
\tableofcontents
\newpage

% ================================================================== 1
\section{Introduction}

\subsection{Contexte}

Un commerçant en ligne doit commander ses produits avant de les vendre. S'il commande trop,
il immobilise de la trésorerie et de l'espace ; s'il commande trop peu, il perd des ventes.
Une prévision de la quantité vendue la semaine suivante, produit par produit, éclaire
directement cette décision de réassort. Le périmètre retenu est volontairement étroit : les
\textbf{50 produits} qui concentrent l'essentiel des volumes, le marché du
\textbf{Royaume-Uni}, majoritaire dans les données, et un horizon d'\textbf{une semaine ISO}.

\subsection{Problématique}

Entraîner un modèle de prévision une fois, dans un notebook, est simple. Le difficile
commence lorsque ce modèle doit vivre : de nouvelles données arrivent chaque semaine et
peuvent être erronées, un nouveau modèle peut être moins bon que l'ancien, un service peut
tomber, et il faut pouvoir expliquer après coup quelle version a produit quelle prévision.
La question posée par ce projet est donc la suivante :

\begin{encadre}[title=Problématique]
Comment faire tourner un modèle de prévision réentraîné chaque semaine de façon
\textbf{fiable} (des données contrôlées, des incidents absorbés), \textbf{traçable} (chaque
prévision rattachée à une version, chaque version à ses données) et
\textbf{exploitable} (un service mesuré, supervisé et dimensionné) ?
\end{encadre}

Le cahier des charges précise que ni la performance ni la complexité du modèle ne sont
évaluées. Le modèle est donc volontairement simple : il n'est qu'un passager, et l'essentiel
du travail porte sur la chaîne qui le fabrique, le contrôle, le met en service, le surveille
et permet de revenir en arrière.

\subsection{Démarche et plan}

Le projet a été construit par couches successives, chacune vérifiée avant de passer à la
suivante : données et qualité, entraînement et registre, orchestration, service,
persistance, supervision, puis tests de charge et démonstration. Les incidents rencontrés
en cours de route ont été traités comme des résultats à part entière : chacun a donné lieu à
une correction et, le plus souvent, à un indicateur qui permet de le détecter.

Le rapport suit cette progression. La section~\ref{sec:donnees} présente les données et leur
rejeu, la section~\ref{sec:archi} l'architecture et les choix techniques. Les sections
\ref{sec:code} à \ref{sec:supervision} décrivent chaque brique de la chaîne, la
section~\ref{sec:charge} les tests de charge et le dimensionnement, et la
section~\ref{sec:demo} la démonstration de bout en bout. La section~\ref{sec:bilan} croise
les critères d'évaluation avec les preuves apportées, avant les limites et la conclusion.

% ================================================================== 2
\section{Les données}\label{sec:donnees}

\subsection{Provenance et conditions d'utilisation}

Le jeu de données \emph{Online Retail} rassemble les transactions d'un commerce en ligne
britannique d'articles cadeaux. Il est publié par le dépôt de l'Université de Californie à
Irvine sous une licence ouverte qui autorise la réutilisation à condition de citer la
source (tableau~\ref{tab:fiche}).

\begin{table}[htbp]\centering\small
\begin{tabularx}{\linewidth}{>{\bfseries}l X}
\toprule
Jeu de données & \emph{Online Retail}, UCI Machine Learning Repository\\
Page & \url{https://archive.ics.uci.edu/dataset/352/online+retail}\\
DOI & 10.24432/C5BW33\\
Licence & Creative Commons Attribution 4.0 (CC BY 4.0) : réutilisation et modification
autorisées, y compris commerciales, sous réserve de citer la source\\
Producteur & Daqing Chen, London South Bank University\\
Contenu & 541\,909 lignes de facture, du 01/12/2010 au 09/12/2011\\
\bottomrule
\end{tabularx}
\caption{Fiche d'identité du jeu de données.}\label{tab:fiche}
\end{table}

Chaque ligne décrit un article d'une facture : numéro de facture (préfixé par \code{C} pour
une annulation), code et libellé du produit, quantité (négative pour un retour), date, prix
unitaire en livres sterling, identifiant client et pays. La seule information relative aux
personnes est \code{CustomerID}, un identifiant numérique déjà pseudonymisé à la source ;
il n'entre pas dans le modèle, puisque la prévision est agrégée par produit, et n'est jamais
exposé par l'API.

Les données ne sont pas versionnées dans Git, à la fois pour leur taille et pour respecter
leur mode de diffusion. La commande \code{retail data} les télécharge, vérifie l'empreinte
SHA256 de l'archive, figée dans \code{configs/data.yaml}, et refuse toute archive
différente : deux personnes qui reproduisent le projet travaillent donc nécessairement sur
les mêmes octets.

\subsection{Simuler l'arrivée de nouvelles données}

Le jeu de données est historique, alors qu'une chaîne opérationnelle doit absorber des
données qui arrivent au fil de l'eau. Il est donc \textbf{rejoué par lots hebdomadaires} :
chaque lot regroupe les factures d'une semaine ISO, du lundi au dimanche, et l'orchestrateur
le « livre » comme si la semaine venait de se terminer. Les première et dernière semaines,
incomplètes, sont écartées ; il reste \textbf{51 lots}, de \code{2010-W49} à
\code{2011-W48}. La semaine \code{2010-W52}, sans ventes en raison de la fermeture de Noël,
est enregistrée comme manquante dans le manifeste plutôt que masquée, et chaque lot y porte
une empreinte de son contenu.

Ce rejeu a une conséquence utile : le lot de la semaine N+1 apporte les ventes réelles de la
semaine prévue lors du lot N. C'est ce qui permet de mesurer l'erreur du modèle \emph{en
production}, et non seulement à l'entraînement (section~\ref{sec:persistance}).

Une précaution évite enfin toute fuite d'information du futur vers le passé : la liste des
50 produits suivis est calculée sur les 8 premières semaines seulement, puis figée.

% ================================================================== 3
\section{Architecture et choix techniques}\label{sec:archi}

\subsection{Vue d'ensemble}

La figure~\ref{fig:archi} montre la circulation des données et des artefacts. Le flux
principal se lit de gauche à droite : un lot est contrôlé, transformé en variables, sert à
entraîner un modèle qui est enregistré, puis, s'il est retenu, servi par l'API. Les traces
d'usage reviennent vers la base de données, et la supervision observe l'ensemble.

\begin{figure}[htbp]\centering
\begin{tikzpicture}[
  font=\small,
  bloc/.style={draw=bleu,fill=bleuclair,rounded corners=3pt,minimum width=2.25cm,
               minimum height=0.95cm,align=center,line width=0.6pt},
  stock/.style={draw=gris,fill=fondcode,cylinder,shape border rotate=90,aspect=0.18,
               minimum width=2.2cm,minimum height=1.05cm,align=center},
  ext/.style={draw=orange,fill=orange!10,rounded corners=3pt,minimum width=2.1cm,
               minimum height=0.9cm,align=center},
  fl/.style={-{Stealth[length=2.4mm]},line width=0.7pt,color=bleu},
  fb/.style={-{Stealth[length=2.2mm]},line width=0.6pt,color=vert,dashed},
  node distance=0.55cm and 0.5cm]

  \node[stock] (uci) {Archive UCI\\\scriptsize SHA256 figé};
  \node[bloc,right=of uci] (lots) {Lots ISO\\\scriptsize 51 semaines};
  \node[bloc,right=of lots] (qual) {Qualité\\\scriptsize Pandera};
  \node[bloc,right=of qual] (feat) {Features\\\scriptsize lags, moyenne};
  \node[bloc,right=of feat] (train) {Entraînement\\\scriptsize Ridge vs naïf};

  \node[stock,below=1.3cm of train] (mlf) {MLflow\\\scriptsize registre :5001};
  \node[bloc,left=1.1cm of mlf] (champ) {Champion\\\scriptsize export disque};
  \node[bloc,left=1.1cm of champ] (api) {API FastAPI\\\scriptsize :8001};
  \node[stock,left=1.1cm of api] (pg) {PostgreSQL\\\scriptsize :5437};

  \node[ext,below=1.25cm of api] (prom) {Prometheus\\\scriptsize :9091};
  \node[ext,left=of prom] (graf) {Grafana\\\scriptsize :3001};
  \node[ext,right=of prom] (loc) {Locust\\\scriptsize charge};

  \begin{scope}[on background layer]
    \node[draw=bleu!60,dashed,rounded corners=5pt,fit=(lots)(train),inner sep=7pt,
          label={[font=\scriptsize\bfseries,text=bleu,anchor=south east]north east:Prefect : un lot par tick}] {};
  \end{scope}

  \draw[fl] (uci) -- (lots);
  \draw[fl] (lots) -- (qual);
  \draw[fl] (qual) -- node[above,font=\scriptsize]{OK} (feat);
  \draw[fl] (feat) -- (train);
  \draw[fl] (train) -- node[right,font=\scriptsize]{enregistre} (mlf);
  \draw[fl] (mlf) -- node[above,font=\scriptsize]{promeut} (champ);
  \draw[fl] (champ) -- node[above,font=\scriptsize]{sert} (api);
  \draw[fl] (api) -- node[above,font=\scriptsize]{traces} (pg);
  \draw[fb] (qual.south) |- ++(0,-0.55) -| node[pos=0.25,above,font=\scriptsize,text=vert]{feedback auto} (pg.north);
  \draw[fl,color=rouge] (qual.north) -- ++(0,0.85) node[above,font=\scriptsize,text=rouge]{KO : arrêt + alerte};
  \draw[fl,color=orange] (prom) -- node[left,font=\scriptsize]{/metrics} (api);
  \draw[fl,color=orange] (graf) -- (prom);
  \draw[fl,color=orange] (graf.north) -- (pg.south);
  \draw[fl,color=orange] (loc) -- (api.south east);
\end{tikzpicture}
\caption{Architecture et circulation des données et des artefacts.}\label{fig:archi}
\end{figure}

Trois principes ont guidé la conception. D'abord, \textbf{un seul point d'entrée} : toutes
les opérations passent par la commande \code{retail}, ce qui rend chaque étape scriptable,
testable et reproductible sans notebook. Ensuite, \textbf{arrêter plutôt que propager} :
une donnée non conforme stoppe la chaîne au lieu de produire un modèle faux. Enfin,
\textbf{le service avant la trace} : une panne de la base ou du registre dégrade la
traçabilité, mais ne doit jamais interrompre les prédictions.

\subsection{Choix techniques}

Les outils ont été choisis pour leur adéquation à une exécution locale, gratuite et
reproductible, et pour la clarté de ce qu'ils permettent de démontrer. Le
tableau~\ref{tab:choix} résume chaque choix, sa justification et l'alternative écartée.

\begin{table}[htbp]\centering\small
\begin{tabularx}{\linewidth}{>{\bfseries}l l X X}
\toprule
Besoin & Choix & Justification & Alternative écartée\\
\midrule
Configuration & YAML + \code{.env} & règles métier relues en revue de code ; secrets hors Git & variables codées en dur\\
Qualité & Pandera & schémas déclaratifs en Python, contrôles nommés, rapport par lot & Great Expectations, plus lourd à configurer pour ce volume\\
Modèle & Ridge & simple, rapide (quelques secondes), comparable à une prévision naïve & modèles d'arbres, sans intérêt ici puisque le score n'est pas évalué\\
Registre & MLflow & runs, versions, alias \code{champion}/\code{challenger}, tags & simple dossier de fichiers, sans historique ni décision tracée\\
Orchestration & Prefect & flux en Python, relances conditionnelles, planification, interface & Airflow, plus lourd à déployer en local\\
Service & FastAPI & validation des entrées par Pydantic, documentation automatique & Flask, sans validation native\\
Traces & PostgreSQL & partagé par l'API, le pipeline et Grafana, robuste en écriture concurrente & SQLite, inadapté à plusieurs processus\\
Supervision & Prometheus + Grafana & standard, alertes déclaratives, tableau de bord versionné & journaux seuls, sans seuil ni alerte\\
Charge & Locust + Docker & scénarios en Python, mêmes paliers sur plusieurs limites CPU et mémoire & tests manuels non reproductibles\\
\bottomrule
\end{tabularx}
\caption{Choix techniques, justification et alternative écartée.}\label{tab:choix}
\end{table}

\paragraph{Déploiement local.} Les services d'infrastructure (PostgreSQL, MLflow,
Prometheus, Grafana) tournent dans Docker Compose ; l'API et le pipeline tournent sur le
poste. Les ports sont fixés dans \code{.env.example} : PostgreSQL 5437, MLflow 5001,
Prefect 4201, API 8001, Grafana 3001 et Prometheus 9091.

% ================================================================== 4
\section{Code et reproductibilité}\label{sec:code}

\subsection{Organisation du code}

Le code forme un paquet Python installable, \code{retail}, découpé par responsabilité
(tableau~\ref{tab:modules}). Chaque module peut être testé isolément, et les dépendances
entre modules suivent le sens du flux : la qualité ne connaît pas l'API, l'API ne connaît pas
l'entraînement.

\begin{table}[htbp]\centering\small
\begin{tabular}{ll}
\toprule
Module & Responsabilité\\
\midrule
\code{common/} & configuration (Pydantic), journalisation JSON, adresse MLflow\\
\code{data/} & téléchargement, vérification d'empreinte, découpage en lots\\
\code{quality/} & schémas Pandera, nettoyage tracé, validation bloquante\\
\code{features/} & construction du panel et des variables\\
\code{training/} & entraînement, métriques, comparaison au naïf\\
\code{registry/} & enregistrement, promotion, retour arrière, export\\
\code{orchestration/} & curseur de rejeu, exécutions, alertes, verrou\\
\code{api/} & service FastAPI, état du modèle servi, cache des variables\\
\code{db/} & tables SQLAlchemy, écriture en file, feedback automatique\\
\bottomrule
\end{tabular}
\caption{Modules du paquet \code{retail}.}\label{tab:modules}
\end{table}

\subsection{Configuration, dépendances et tests}

La configuration est séparée en deux sources. Les fichiers \code{configs/config.yaml} et
\code{configs/data.yaml}, versionnés, contiennent tout ce qui relève du métier : seuils de
qualité, paramètres d'entraînement, règle de promotion, source des données et son
empreinte. Le fichier \code{.env}, jamais versionné, contient les adresses et identifiants
des services ; un modèle \code{.env.example} sans secret l'accompagne. L'ensemble est validé
au démarrage par Pydantic : une valeur hors bornes arrête la commande avant tout traitement,
plutôt qu'au milieu d'un entraînement.

Les dépendances sont maîtrisées à deux niveaux : des bornes de compatibilité dans
\code{pyproject.toml}, et des versions exactes verrouillées dans \code{requirements.lock}
pour Python 3.11 à 3.13. La qualité du code est contrôlée par \code{ruff} (règles et
formatage), et 18 fichiers de tests unitaires et d'intégration couvrent la configuration,
l'ingestion, la qualité, l'entraînement, le registre, l'orchestration, l'API, la base, la
supervision, la démonstration et les tests de charge. Une intégration continue GitHub
Actions rejoue le lint et les tests sur trois versions de Python.

\subsection{Vérification depuis un clone neuf}

Une installation qui ne fonctionne que sur le poste de son auteur n'est pas reproductible.
Le dépôt a donc été cloné depuis GitHub dans un environnement vide, puis installé et
exécuté pas à pas, comme le ferait un correcteur (tableau~\ref{tab:clone}). Cette
vérification a d'ailleurs révélé un défaut réel, corrigé depuis : une règle trop large du
\code{.gitignore} excluait le dossier \code{src/retail/data/}, si bien que le dépôt publié ne
démarrait pas.

\begin{table}[htbp]\centering\small
\begin{tabularx}{\linewidth}{X l}
\toprule
Étape & Résultat\\
\midrule
Installation (\code{requirements.lock}, puis \code{pip install -e .}) & réussie\\
\code{ruff check} et \code{ruff format --check} & aucune erreur, 66 fichiers conformes\\
\code{retail check-config} & configuration valide\\
\code{pytest} & \textbf{131 tests réussis}\\
\code{retail data} & empreinte vérifiée, 51 lots produits\\
\code{retail quality}, puis lot invalide de démonstration & lots valides acceptés, lot invalide rejeté\\
\code{retail train}, \code{retail replay}, \code{retail rollback} & chaîne complète exécutée\\
API : \code{/health}, \code{/predict/batch} & modèle servi, 50 prédictions, version renvoyée\\
\bottomrule
\end{tabularx}
\caption{Rejeu de la chaîne depuis un clone neuf du dépôt.}\label{tab:clone}
\end{table}

% ================================================================== 5
\section{Qualité des données}\label{sec:qualite}

\subsection{Des contrôles qui bloquent}

Un modèle entraîné sur des données fausses produit des prévisions fausses, et rien dans ses
métriques ne le signale nécessairement. La qualité est donc contrôlée \emph{avant} tout
prétraitement, et un lot non conforme arrête la chaîne. Chaque lot traverse trois étapes :

\begin{enumerate}
  \item le \textbf{schéma des factures} vérifie les colonnes et leurs types, que toutes les
        dates tombent dans la semaine ISO du lot, que le lot compte au moins 500 lignes et
        que moins de 5\,\% des prix sont nuls ou négatifs ;
  \item le \textbf{nettoyage} écarte les ventes hors Royaume-Uni, les annulations, les
        quantités ou prix nuls, les codes qui ne sont pas des produits (frais de port,
        remises\ldots) et les valeurs aberrantes ; chaque ligne écartée est comptée et
        justifiée dans un rapport JSON, ce qui rend le nettoyage auditable ;
  \item le \textbf{schéma du panel} vérifie le résultat agrégé : une ligne par produit et par
        semaine, des quantités positives ou nulles, aucun doublon.
\end{enumerate}

Si l'une de ces étapes échoue, aucun panel n'est écrit, l'entraînement n'est pas lancé, une
alerte est émise et le curseur du rejeu reste bloqué sur la semaine fautive jusqu'à
correction.

\subsection{Distinguer erreur de données et panne technique}

Une erreur de données et une panne technique n'appellent pas la même réponse. Une base
momentanément injoignable se résout souvent d'elle-même : il est raisonnable de réessayer.
Un lot invalide, lui, le restera : le réessayer ne ferait que retarder l'alerte. Les deux
cas sont donc représentés par des exceptions distinctes, et l'orchestrateur ne relance
automatiquement que les pannes techniques (section~\ref{sec:orchestration}).

\subsection{Mise à l'épreuve}

La commande \code{retail quality --invalid-demo} fabrique un lot volontairement corrompu,
\code{2099-W01}. Lors de la démonstration, il a été rejeté par trois contrôles
(dates hors de la semaine, part de prix non positifs, volume minimum), aucun entraînement
n'a été lancé, et l'alerte correspondante a été reçue par l'API.

% ================================================================== 6
\section{Modèle et entraînement}\label{sec:modele}

Les ventes sont agrégées en un panel « produit $\times$ semaine », complété par des zéros
pour les semaines sans vente. Le modèle s'appuie sur les ventes des 1\textsuperscript{re},
2\textsuperscript{e} et 4\textsuperscript{e} semaines précédentes et leur moyenne mobile sur
quatre semaines, sur le prix moyen de la semaine précédente, sur la semaine de l'année et le
mois pour la saisonnalité, et sur le code produit. Le prix est pris à la semaine
\emph{précédente} : utiliser celui de la semaine prévue reviendrait à exploiter une
information inconnue au moment de la prévision.

L'évaluation respecte l'ordre du temps : les 8 dernières semaines disponibles forment le jeu
de test et ne sont jamais utilisées pour apprendre. Le modèle, une régression Ridge, est
comparé à une \textbf{prévision naïve} qui reconduit les ventes de la semaine précédente.
Deux métriques sont suivies : la \textbf{MAE}, erreur absolue moyenne en unités vendues, qui
sert aux décisions, et la \textbf{WAPE}, erreur rapportée au volume vendu, plus parlante pour
le métier.

Ce choix d'un modèle linéaire est délibéré. La performance n'étant pas évaluée, un modèle
plus complexe n'apporterait rien à la démonstration ; en revanche, un entraînement de
quelques secondes rend possible un rejeu automatique toutes les dix minutes.

% ================================================================== 7
\section{Traçabilité et cycle de vie des modèles}\label{sec:registre}

\subsection{Tracer pour pouvoir expliquer}

Chaque entraînement crée un run MLflow qui conserve tout ce qu'il faut pour l'expliquer ou
le reproduire : la dernière semaine utilisée (\code{until_week}) et les semaines de test,
une empreinte des variables (\code{features_hash}) et de l'archive source
(\code{data_hash}), les paramètres, les métriques du modèle et de la prévision naïve, et les
artefacts, c'est-à-dire le modèle avec sa signature et la liste figée des produits.

\subsection{Une règle de mise en service explicite}

Un nouveau modèle n'est pas automatiquement meilleur que l'ancien. La règle de mise en
service, écrite dans la configuration, procède en deux temps. Un run n'est
\textbf{enregistré} comme version candidate (alias \code{challenger}) que s'il bat la
prévision naïve. Il n'est ensuite \textbf{promu} champion que si sa MAE est inférieure ou
égale à celle du champion en place, \textbf{mesurée sur le même jeu de test}. Ce dernier
point est essentiel : deux modèles évalués sur des semaines différentes ne sont pas
comparables, puisque certaines semaines sont intrinsèquement plus difficiles à prévoir.

La décision est écrite en tag sur chaque version, ce qui permet de la relire après coup.
Pour la version 7, par exemple :
\begin{lstlisting}
decision : promoted: mae 242.9378 <= champion 6 mae 243.0059 on same test
previous_champion : 6
\end{lstlisting}

Un garde-fou complète la règle : si une semaine déjà entraînée est rejouée, aucune nouvelle
version n'est créée. Sans lui, un ré-entraînement sur les mêmes données produisait une copie
identique du modèle, promue grâce au « inférieur ou égal », ce qui brouillait l'historique.

\subsection{Servir sans dépendre du registre}

À chaque promotion, le champion est exporté de façon atomique dans
\code{models/champion/}, avec une copie de secours. Si le registre MLflow devient
injoignable, l'API continue de servir ce champion local et le signale comme \textbf{mode
dégradé}, dans \code{/health} comme dans les métriques.

\subsection{Revenir en arrière}

Si une version promue se révèle décevante en production, la commande \code{retail rollback}
remet en service la version indiquée par son tag \code{previous_champion}, ou une version
choisie avec \code{--to-version}. Le champion est réexporté à chaque fois, et les retours
arrière peuvent s'enchaîner. La démonstration l'a vérifié en conditions réelles : passage de
la version 7 à la version 6, puis restauration de la 7.

% ================================================================== 8
\section{Orchestration}\label{sec:orchestration}

\subsection{Le flux hebdomadaire}

Le flux Prefect \code{weekly_lot} traite \textbf{un lot par exécution}. Il est déclenché
automatiquement toutes les dix minutes, chaque exécution représentant une semaine simulée,
et une seule exécution peut tourner à la fois. Chaque exécution choisit la semaine suivante
à partir du curseur de rejeu, valide le lot, rapproche les ventes réelles des prédictions
faites pour cette semaine, entraîne puis enregistre éventuellement un nouveau modèle, et
demande enfin à l'API la prévision de la semaine suivante, qui sera comparée aux ventes
réelles au lot suivant. La boucle se ferme ainsi d'elle-même.

\subsection{Gestion des échecs}

L'enjeu de l'orchestration n'est pas l'enchaînement nominal, simple, mais le comportement
face aux échecs (tableau~\ref{tab:echecs}).

\begin{table}[htbp]\centering\small
\begin{tabularx}{\linewidth}{l X}
\toprule
Situation & Comportement\\
\midrule
Lot non conforme & arrêt, alerte, curseur bloqué sur la semaine, aucune relance\\
Panne technique pendant l'entraînement & 2 relances espacées de 10 s, puis alerte\\
Interruption au milieu d'un enregistrement & reprise de la version restée « en cours », sans ré-entraîner\\
Tous les lots déjà rejoués & exécution terminée normalement, sans alerte\\
API arrêtée au moment de la prévision & le lot reste validé ; l'échec de la prévision est noté\\
Rejeu manuel d'une semaine passée & le curseur ne recule pas\\
\bottomrule
\end{tabularx}
\caption{Comportement de l'orchestration face aux échecs.}\label{tab:echecs}
\end{table}

\subsection{Suivi des exécutions}

Chaque exécution est enregistrée dans un journal (\code{runs.jsonl}) et dans la table
\code{pipeline_runs}, avec la semaine, le statut, la durée, la version et l'éventuelle
erreur ; Prefect conserve en parallèle l'historique détaillé de chaque tâche. Les alertes
sont envoyées à l'API, qui les stocke ; une copie locale est écrite si l'envoi échoue. Un
verrou de fichier empêche enfin deux rejeux simultanés de se marcher dessus.

% ================================================================== 9
\section{API de prédiction}\label{sec:api}

L'API expose les prédictions en ligne et par lot, ainsi que l'état du service
(tableau~\ref{tab:routes}).

\begin{table}[htbp]\centering\small
\begin{tabularx}{\linewidth}{l l X}
\toprule
Méthode & Route & Rôle\\
\midrule
POST & \code{/predict} & prédiction pour un produit et une semaine\\
POST & \code{/predict/batch} & prédictions pour les 50 produits d'une semaine\\
POST & \code{/predict/features} & prédiction à partir de variables fournies\\
POST & \code{/feedback} & retour terrain (quantité réelle) sur une prédiction\\
GET & \code{/health} & état complet : version, source, mode dégradé, base\\
GET & \code{/health/live} & vivacité, sans appel aux dépendances\\
GET & \code{/model} & version servie, produits, WAPE de production\\
POST & \code{/model/reload} & rechargement du champion\\
GET/POST & \code{/alerts} & réception et consultation des alertes du pipeline\\
GET & \code{/metrics} & métriques Prometheus\\
\bottomrule
\end{tabularx}
\caption{Routes de l'API.}\label{tab:routes}
\end{table}

Les entrées sont validées par Pydantic avant tout calcul : format de semaine
\code{AAAA-Www}, champs obligatoires, et horizon limité à la semaine qui suit la dernière
semaine connue, puisqu'au-delà les variables ne peuvent pas être calculées. Les erreurs sont
explicites : \textbf{404} pour un produit inconnu ou une prédiction introuvable,
\textbf{422} pour une entrée invalide ou hors horizon.

Chaque réponse identifie la version du modèle qui l'a produite, dans le corps
(\code{model_version}) et dans l'en-tête \code{X-Model-Version}, et porte un identifiant
unique de prédiction qui sert à rattacher le retour terrain. Le modèle servi est remplacé de
façon atomique : une requête en cours ne voit jamais un état à moitié mis à jour. Chaque
worker vérifie toutes les 30 secondes si le champion a changé.

% ================================================================== 10
\section{Persistance et retours terrain}\label{sec:persistance}

\subsection{Conserver les traces d'usage}

Quatre tables PostgreSQL conservent l'activité : les \textbf{prédictions} (produit,
semaine, quantité prédite, version, durée, statut et erreur éventuelle), les
\textbf{feedbacks} (quantité réelle, quantité prédite, erreur, version, origine), les
\textbf{exécutions du pipeline} et les \textbf{alertes}.

L'écriture des prédictions suit le principe « le service avant la trace » : elles passent
par une file bornée vidée en arrière-plan, par lots. Si la base est lente ou coupée, l'API
continue de répondre, et les écritures perdues sont \emph{comptées} par une métrique au lieu
de disparaître en silence. Les exécutions et alertes du pipeline, peu nombreuses, sont
écrites directement. Le trafic des tests de charge, enfin, est marqué par un en-tête dédié
et exclu des retours terrain, pour ne pas fausser les indicateurs de production.

\subsection{Mesurer l'erreur en production}

Quand le lot d'une semaine est validé, ses ventes réelles ferment automatiquement les
prédictions faites pour cette semaine. On obtient ainsi une \textbf{WAPE de production par
version}, visible dans \code{/model} et dans Grafana. Elle complète utilement la MAE
d'entraînement : la première dit comment le modèle s'est comporté sur les semaines de test,
la seconde comment il se comporte réellement une fois en service. Sur les semaines rejouées,
elle se situe entre 62\,\% et 71\,\% selon la version, un niveau cohérent avec la simplicité
du modèle et la volatilité des ventes hebdomadaires de produits cadeaux. L'API accepte aussi
des retours saisis manuellement, rattachés à la prédiction d'origine.

% ================================================================== 11
\section{Supervision}\label{sec:supervision}

\subsection{Le tableau de bord}

Le tableau de bord Grafana \code{retail-mlops} est provisionné automatiquement : il fait
partie du dépôt et s'installe avec les services. Il répond à quatre questions d'exploitation.
\emph{Le service est-il utilisé ?} : requêtes par seconde et version servie.
\emph{Répond-il assez vite ?} : latences p50, p95 et p99, avec un seuil à 200 ms.
\emph{Est-il en bonne santé ?} : taux d'erreurs 4xx et 5xx, mode dégradé, état de la base,
écritures abandonnées et frise des alertes. \emph{Le modèle et le pipeline se
comportent-ils bien ?} : exécutions du pipeline, alertes du pipeline, WAPE de production par
version et par semaine, volumes de traces.

La figure~\ref{fig:grafana.png} montre ce tableau de bord pendant la démonstration : la
coupure volontaire de PostgreSQL y est visible sous trois angles, par les écritures
abandonnées, par l'alerte qui passe d'« en attente » à « déclenchée », et par l'absence
d'erreurs 5xx, preuve que le service a continué de répondre.

\capture[0.97\linewidth]{grafana.png}{Tableau de bord Grafana sur la fenêtre de la démonstration (10 novembre 2026, 03:40--03:50). Pendant la coupure volontaire de PostgreSQL vers 03:45, les écritures abandonnées montent et l'alerte \code{RetailDbDown} passe d'« en attente » à « déclenchée ». En dessous : WAPE de production par version et par semaine, exécutions du pipeline (durée, version), alertes du pipeline et volumes de traces.}

\subsection{Les alertes}

Quatre règles Prometheus couvrent les situations qui exigent une intervention
(tableau~\ref{tab:alertes}). Leurs délais sont choisis pour ignorer un incident fugace tout
en signalant rapidement un problème durable.

\begin{table}[htbp]\centering\small
\begin{tabularx}{\linewidth}{l l X}
\toprule
Alerte & Condition & Signification\\
\midrule
\code{RetailDegraded} & mode dégradé pendant 1 min & le registre est injoignable, le champion local est servi\\
\code{RetailHigh5xx} & plus de 5\,\% de 5xx pendant 1 min & le service renvoie des erreurs\\
\code{RetailApiDown} & métriques injoignables pendant 30 s & l'API est arrêtée\\
\code{RetailDbDown} & base injoignable pendant 15 s & les traces ne sont plus écrites\\
\bottomrule
\end{tabularx}
\caption{Règles d'alerte Prometheus.}\label{tab:alertes}
\end{table}

% ================================================================== 12
\section{Tests de charge et dimensionnement}\label{sec:charge}

\subsection{Protocole}

L'objectif est de répondre à une question concrète : combien de ressources faut-il pour
tenir une charge donnée avec une latence acceptable ? Le protocole, piloté par
\code{load_tests/run.py}, simule avec Locust un usage réaliste : 70\,\% de prédictions
unitaires, 20\,\% de prédictions par lot et 10\,\% de sondes de santé. La charge monte par
paliers de 2, 5, 10, 20, 40, 80, 160 puis 320 utilisateurs simultanés, chacun pendant
45 secondes. Trois configurations sont comparées, chacune dans un conteneur limité en CPU et
en mémoire : 1 worker sur 0,5 vCPU, 2 workers sur 1 vCPU, 4 workers sur 2 vCPU. La
\textbf{capacité} d'une configuration est définie comme le dernier palier où la latence p95
reste sous 200 ms sans aucune erreur.

\subsection{Une première série qui a révélé deux défauts}

La première série de mesures n'atteignait le critère dans aucune configuration : le débit
plafonnait entre 9,5 et 38,8 requêtes par seconde, avec un p95 toujours au-delà d'une
seconde. Son analyse a mis au jour deux défauts. À chaque requête, l'API relisait la liste
des fichiers du panel sur un dossier monté depuis Windows, ce qui coûtait environ 400 ms. Et
79\,\% des traces de prédiction étaient perdues sous charge alors que le compteur d'abandons
restait à zéro : la perte était silencieuse. Après correction (lecture mise en cache 5 s,
insertion SQL par lot, compteur réparé), le débit maximal a été multiplié par 2 à 3,4 et
pratiquement toutes les traces sont enregistrées. Cet épisode illustre l'intérêt d'un
protocole de charge : il a fait apparaître un problème de traçabilité qu'aucun test
fonctionnel n'aurait détecté.

\subsection{Résultats}

\begin{figure}[htbp]\centering
\begin{tikzpicture}
\begin{axis}[
  width=0.49\linewidth, height=6cm,
  xmode=log, log basis x=2, xtick={2,5,10,20,40,80,160,320},
  xticklabels={2,5,10,20,40,80,160,320},
  xlabel={Utilisateurs simultanés}, ylabel={Débit (req/s)},
  ymin=0, ymax=90, ytick={0,20,40,60,80}, grid=major, grid style={gray!25},
  legend style={font=\scriptsize,at={(0.5,1.03)},anchor=south,legend columns=2,draw=none,fill=none},
  mark size=1.6pt, thick]
\addplot[color=orange,mark=*] coordinates {(2,12.4)(5,29.8)(10,32.4)(20,29.3)(40,23.4)(80,21.7)(160,22.3)(320,21.5)};
\addplot[color=bleu,mark=square*] coordinates {(2,12.9)(5,31.0)(10,55.8)(20,49.9)(40,40.0)(80,33.3)(160,33.5)(320,34.7)};
\addplot[color=vert,mark=triangle*] coordinates {(2,12.3)(5,31.3)(10,60.2)(20,81.6)(40,77.0)(80,80.8)(160,68.4)(320,58.0)};
\legend{{1w / 0,5 vCPU},{2w / 1 vCPU},{4w / 2 vCPU}}
\end{axis}
\end{tikzpicture}\hfill
\begin{tikzpicture}
\begin{axis}[
  width=0.49\linewidth, height=6cm,
  xmode=log, log basis x=2, xtick={2,5,10,20,40,80,160,320},
  xticklabels={2,5,10,20,40,80,160,320},
  ymode=log, xlabel={Utilisateurs simultanés}, ylabel={p95 (ms)},
  grid=major, grid style={gray!25},
  legend style={font=\scriptsize,at={(0.5,1.03)},anchor=south,legend columns=2,draw=none,fill=none},
  mark size=1.6pt, thick]
\addplot[color=orange,mark=*] coordinates {(2,51)(5,120)(10,310)(20,1500)(40,3800)(80,6000)(160,6800)(320,17000)};
\addplot[color=bleu,mark=square*] coordinates {(2,27)(5,49)(10,110)(20,720)(40,1900)(80,4100)(160,9800)(320,16000)};
\addplot[color=vert,mark=triangle*] coordinates {(2,34)(5,49)(10,60)(20,360)(40,1200)(80,2100)(160,4900)(320,11000)};
\addplot[color=rouge,dashed,domain=2:320,samples=2] {200};
\legend{1w, 2w, 4w, seuil 200 ms}
\end{axis}
\end{tikzpicture}
\caption{Débit (à gauche) et latence p95 (à droite) par palier, après optimisation
(échelles logarithmiques pour les utilisateurs et le p95).}\label{fig:charge}
\end{figure}

\begin{table}[htbp]\centering\small
\begin{tabular}{lcccccc}
\toprule
Config. & Palier & \textbf{Capacité} & p95 & RAM & Débit max & /vCPU\\
\midrule
1w / 0,5 vCPU & 5 users & \textbf{29,8 req/s} & 120 ms & 305 Mo & 32,4 req/s & 65\\
2w / 1 vCPU & 10 users & \textbf{55,8 req/s} & 110 ms & 550 Mo & 55,8 req/s & 56\\
4w / 2 vCPU & 10 users & \textbf{$\geq$ 60,2 req/s} & 60 ms & 1,05 Go & 81,6 req/s & 41\\
\bottomrule
\end{tabular}
\caption{Capacité mesurée (p95 < 200 ms, 0\,\% d'erreur) et débit maximal par configuration.
La colonne « /vCPU » rapporte le débit maximal au nombre de vCPU.}\label{tab:capacite}
\end{table}

La figure~\ref{fig:charge} et le tableau~\ref{tab:capacite} conduisent à trois constats.
Premièrement, le débit croît avec les ressources, mais de moins en moins vite : 65, puis 56,
puis 41 requêtes par seconde et par vCPU. Au-delà de 2 vCPU, une autre limite apparaît, que
le banc de test ne permet pas de départager (générateur de charge, écriture en base ou poste
lui-même). Deuxièmement, une fois la saturation atteinte, le débit ne plafonne pas : il
\emph{baisse}, pendant que la latence explose, parce que les requêtes s'accumulent en file
d'attente. Il faut donc dimensionner avec une marge, et non au plus juste. Troisièmement, la
mémoire augmente d'environ 250 Mo par worker, ce qui reste modeste. Pour 4 workers, le
palier de 10 utilisateurs n'est pas saturé et le suivant dépasse 200 ms : la capacité réelle
se situe entre 60 et 82 requêtes par seconde.

\subsection{Recommandation}

\begin{encadre}[title=Dimensionnement recommandé]
\textbf{Unité : 2 workers, 1 vCPU, 1 Go de mémoire.} Elle tient \textbf{55,8 req/s} avec un
p95 de 110 ms et aucune erreur, soit la meilleure capacité garantie par vCPU mesurée. En
gardant 30\,\% de marge, on compte environ 39 req/s par instance :
\[
  \text{instances} = \left\lceil \frac{\text{pointe (req/s)}}{39} \right\rceil
\]
Exemples : 30 req/s $\rightarrow$ 1 instance ; 100 req/s $\rightarrow$ 3 instances ;
200 req/s $\rightarrow$ 6 instances derrière un répartiteur de charge.
\end{encadre}

Ajouter des instances de cette taille plutôt que grossir une instance unique est préférable
pour deux raisons : le rendement par vCPU décroît au-delà de 1 vCPU, et plusieurs instances
apportent en prime une tolérance à la panne de l'une d'elles.

% ================================================================== 13
\section{Démonstration de bout en bout}\label{sec:demo}

\subsection{Scénario}

L'intégration doit être démontrée, et pas seulement décrite. La commande \code{retail demo}
enchaîne automatiquement un fonctionnement nominal et des incidents contrôlés, et vérifie à
chaque étape le résultat attendu. Le tableau~\ref{tab:demo} présente les résultats observés
lors de l'exécution du 10 novembre 2026, dont la sortie complète est reproduite à la
figure~\ref{fig:demo.png}.

\begin{table}[htbp]\centering\small
\begin{tabularx}{\linewidth}{c l X}
\toprule
& Scène & Résultat observé\\
\midrule
0 & Vérifications & API, MLflow, PostgreSQL et Prometheus joignables ; version 7 servie depuis MLflow ; 4 règles d'alerte chargées\\
1 & Fonctionnement nominal & 50 prédictions pour la semaine W41 ; livraison du lot W41 ; 50 feedbacks automatiques ; WAPE de production de 0,70\\
2 & Lot invalide & \code{2099-W01} rejeté par 3 contrôles ; aucun entraînement ; alerte reçue par l'API\\
3 & Base de données coupée & \code{/predict} répond toujours ; \code{db_ok} passe à faux ; \code{RetailDbDown} déclenchée ; retour à la normale après redémarrage\\
4 & Retour arrière & champion 7 $\rightarrow$ 6, puis restauration de la version 7\\
5 & Bilan & récapitulatif des endroits où vérifier chaque preuve\\
\bottomrule
\end{tabularx}
\caption{Scènes de la démonstration et résultats observés.}\label{tab:demo}
\end{table}

\subsection{Preuves}

Les captures suivantes complètent celle du tableau de bord (figure~\ref{fig:grafana.png}).
Le registre MLflow (figures~\ref{fig:mlflow.png} et~\ref{fig:mlflow_v7.png}) montre les
versions successives, l'alias qui désigne la version servie et la décision de promotion
écrite sur chacune. Prefect (figures~\ref{fig:prefect.png} et~\ref{fig:prefect_runs.png})
montre une exécution déclenchée automatiquement et l'historique des exécutions, y compris
les échecs, sur lesquels revient la section~\ref{sec:limites}.

\capture[0.97\linewidth]{mlflow.png}{Registre MLflow : versions 7 et 6, alias \code{champion} et \code{challenger} sur la version 7, tags de traçabilité (\code{until_week}, \code{features_hash}, \code{previous_champion}, \code{decision}\ldots).}
\capture[0.97\linewidth]{mlflow_v7.png}{Fiche de la version 7 : la décision de promotion est écrite en tag (« MAE 242,94 $\leq$ champion 6 MAE 243,01 sur le même jeu de test »).}
\capture[0.97\linewidth]{prefect.png}{Prefect : exécution \code{venerable-vole}, déclenchée automatiquement par la planification (\emph{auto-scheduled}) : validation du lot 2011-W47, puis entraînement et enregistrement de la version 6, terminée en 10 s. Le retard affiché correspond au rattrapage des exécutions planifiées après le redémarrage du service.}
\capture[0.97\linewidth]{prefect_runs.png}{Historique Prefect sur une semaine. Vert : exécutions réussies ; jaune : exécutions planifiées restées en retard pendant l'arrêt du service ; rouge, orange et gris : échecs, plantage et annulation du 9 novembre, causés par la saturation mémoire du poste et un serveur MLflow figé (section~\ref{sec:limites}).}
\capture[0.95\linewidth]{demo.png}{Sortie complète de \code{retail demo --reset-cursor} (10 novembre 2026) : les six scènes sont validées. Seuls les accents, déformés par la console Windows, ont été rétablis.}

\subsection{Incidents rencontrés et corrigés}

La mise au point a fait apparaître de véritables incidents d'exploitation. Plutôt que de
les masquer, ils ont été documentés et corrigés ; les plus significatifs sont résumés
dans le tableau~\ref{tab:incidents}, et la liste complète figure dans
\code{docs/SYNTHESE_OPS.md}. Leur point commun est instructif : la plupart étaient
\emph{silencieux} tant qu'aucun indicateur ne les mesurait.

\begin{table}[htbp]\centering\small
\begin{tabularx}{\linewidth}{X X}
\toprule
Incident & Correction\\
\midrule
Workers Uvicorn tués au démarrage sous contrainte CPU (délai de santé de 5 s dépassé) &
délai relevé à 60 s, imports allégés\\
Lecture du dossier panel à chaque requête, environ 400 ms sur un volume monté &
cache de 5 s : débit multiplié par 2 à 3\\
79\,\% des traces perdues sous charge, compteur à zéro &
insertion par lot, abandons comptés\\
Métriques vides dans le conteneur (dossier multiprocessus sur un volume Windows) &
dossier en mémoire (\code{tmpfs})\\
Deux processus sur le port 8001, Prometheus observait le mauvais &
procédure : arrêter le conteneur avant l'API du poste\\
Exécutions du pipeline absentes de la base après \code{retail replay} &
écriture directe au lieu d'un fil d'arrière-plan\\
Versions identiques créées en rejouant une même semaine &
réutilisation de la version existante\\
Alerte de 10 s invisible dans Grafana (pas de requête de 15 s) &
pas aligné sur la collecte Prometheus (5 s)\\
Mémoire de Windows épuisée par la machine virtuelle de Docker &
mémoire de Docker limitée, \code{OPENBLAS_NUM_THREADS=1}\\
\bottomrule
\end{tabularx}
\caption{Incidents d'exploitation et corrections.}\label{tab:incidents}
\end{table}

% ================================================================== 14
\section{Bilan au regard des critères d'évaluation}\label{sec:bilan}

Le tableau~\ref{tab:bilan} met en regard chaque critère du cahier des charges et la preuve
correspondante dans le projet.

\begin{table}[htbp]\centering\small
\begin{tabularx}{\linewidth}{>{\bfseries}p{3.2cm} X}
\toprule
Critère & Ce qui est démontré\\
\midrule
Qualité du code (20\,\%) & paquet modulaire, configuration validée au démarrage, 131 tests, lint \code{ruff}, CI sur trois versions de Python (section~\ref{sec:code})\\
Reproductibilité (10\,\%) & dépendances verrouillées, données vérifiées par empreinte, chaîne rejouée avec succès depuis un clone neuf (tableau~\ref{tab:clone})\\
Intégration MLOps (20\,\%) & lot $\rightarrow$ qualité $\rightarrow$ entraînement $\rightarrow$ registre $\rightarrow$ API $\rightarrow$ prévision $\rightarrow$ feedback, automatisé par Prefect ; 7 versions produites (sections~\ref{sec:qualite} à~\ref{sec:orchestration})\\
Fiabilité et traçabilité (20\,\%) & contrôles bloquants, relances limitées aux pannes, règle de promotion tracée, retour arrière, mode dégradé, incidents documentés (sections~\ref{sec:registre} et~\ref{sec:demo})\\
Exploitation du service (20\,\%) & API validée et versionnée, traces et feedbacks persistés, WAPE de production, tableau de bord et alertes (sections~\ref{sec:api} à~\ref{sec:supervision})\\
Dimensionnement (10\,\%) & protocole par paliers sur trois configurations, critère explicite, recommandation chiffrée et justifiée (section~\ref{sec:charge})\\
\bottomrule
\end{tabularx}
\caption{Correspondance entre les critères d'évaluation et les preuves apportées.}\label{tab:bilan}
\end{table}

% ================================================================== 15
\section{Limites connues}\label{sec:limites}

Plusieurs limites subsistent ; elles sont connues et n'affectent pas la démonstration, mais
il serait nécessaire de les traiter avant un usage réel.

\begin{itemize}
  \item \textbf{Banc de charge sur un seul poste.} L'API, le générateur de charge et les
        services de supervision se partagent la même machine ; la limite observée au-delà
        de 2 vCPU ne peut pas être attribuée avec certitude.
  \item \textbf{Latences mesurées côté serveur.} Elles commencent à la prise en charge de la
        requête et excluent l'attente qui la précède ; vue du client, la latence sous charge
        est plus élevée.
  \item \textbf{Ressources du poste.} Le 9 novembre, la machine virtuelle de Docker a réservé
        plus de 12 Go ; la mémoire engagée de Windows s'est épuisée, des exécutions Prefect
        ont échoué et MLflow s'est figé. La mémoire de Docker est désormais limitée, mais
        l'épisode rappelle que l'environnement d'exécution fait partie du système.
  \item \textbf{Prévision faite par la version précédente.} Juste après une promotion, l'API
        peut servir encore l'ancien champion pendant au plus 30 secondes ; la prévision de
        la semaine suivante est alors attribuée, à juste titre, à cette version.
  \item \textbf{Feedback immédiat.} Un retour envoyé dans la seconde qui suit la prédiction
        peut être refusé, le temps que la prédiction soit écrite en base.
  \item \textbf{Alertes enregistrées deux fois.} Une alerte du pipeline est stockée à la fois
        par l'API qui la reçoit et par le pipeline qui l'émet.
  \item \textbf{Empreinte des données.} Le tag \code{data_hash} identifie l'archive complète ;
        la tranche réellement utilisée est donnée par \code{until_week} et
        \code{features_hash}.
  \item \textbf{Exécution locale.} Il n'y a ni haute disponibilité ni déploiement distant,
        conformément au cadre du projet.
\end{itemize}

% ================================================================== 16
\section{Conclusion}

Ce projet répond à la problématique posée en introduction par une chaîne complète et
vérifiée. Elle est \textbf{fiable} : les données sont contrôlées avant tout entraînement, les
erreurs de données sont distinguées des pannes, et une coupure de la base est absorbée sans
interrompre le service. Elle est \textbf{traçable} : chaque version est reliée à ses données
et à la décision qui l'a mise en service, chaque prévision à la version qui l'a produite, et
le retour arrière tient en une commande. Elle est \textbf{exploitable} : le service est
mesuré, supervisé, alerté, et son dimensionnement repose sur des mesures plutôt que sur une
estimation.

La démarche a été autant itérative que constructive. Les tests de charge ont révélé un goulot
d'étranglement et une perte silencieuse de traces ; la vérification depuis un clone neuf a
révélé un dépôt incomplet ; l'exploitation sur plusieurs jours a révélé une saturation
mémoire. Chacun de ces problèmes a été corrigé puis rendu observable. C'est sans doute
l'enseignement principal du projet : en exploitation, ce qui n'est pas mesuré ne peut pas
être garanti.

\paragraph{Pistes d'amélioration.} Mesurer la charge avec un générateur placé sur une autre
machine ; ajouter un palier intermédiaire pour préciser la capacité à 4 workers ; protéger
l'API contre la surcharge par une limite de concurrence ; recharger le champion dans l'API
dès la promotion ; calculer l'empreinte des données sur la seule tranche d'entraînement.

% ================================================================== annexes
\appendix
\section{Reproduire le projet}\label{ann:commandes}

\begin{lstlisting}
# Installation
python -m venv .venv
pip install -r requirements.lock
pip install -e . --no-deps
cp .env.example .env                 # Windows : copy .env.example .env
retail check-config

# Services
docker compose up -d postgres mlflow prometheus grafana
retail db-init

# Données, qualité, premiers modèles
retail data
retail quality
retail train --until-week 2011-W38
retail train --until-week 2011-W40

# Service et démonstration (API dans un autre terminal)
retail api --host 0.0.0.0
retail demo --reset-cursor

# Orchestration
prefect server start --host 127.0.0.1 --port 4201
python flows/weekly_lot.py

# Tests et charge
pytest
python load_tests/run.py
\end{lstlisting}

\section{Références}

\begin{itemize}
  \item Chen, D. (2015). \emph{Online Retail} [Dataset]. UCI Machine Learning Repository.
        \url{https://doi.org/10.24432/C5BW33}
  \item Chen, D., Sain, S. L., Guo, K. (2012). Data mining for the online retail industry:
        A case study of RFM model-based customer segmentation using data mining.
        \emph{Journal of Database Marketing and Customer Strategy Management}, 19(3), 197--208.
  \item Documentations : MLflow, Prefect, FastAPI, Pandera, Prometheus, Grafana, Locust.
\end{itemize}

\end{document}
